\subsection{Part A – Knowledge check}

\indent 

\noindent\textbf{Define IT Governance in one sentence.}


IT Governance is the framework that ensures IT investments and operations are aligned with business goals, deliver value, and manage risks effectively

\bigskip

\noindent\textbf{What are the four domains of COBIT?}


Planning and Organisation, Acquisition and Implementation, Delivery and Support, and Monitoring and Evaluation


\bigskip

\noindent\textbf{List two differences between ITIL v3 and ITIL v4.}


ITIL v3 is structured around the Service Lifecycle, while ITIL v4 is built on the Service Value System (SVS) and Service Value Chain. ITIL v3 focuses mainly on processes, whereas ITIL v4 integrates Agile, DevOps, and digital transformation practices for flexibility.

\bigskip

\noindent\textbf{Name three ITIL v4 guiding principles.}


Focus on Value, Start Where You Are, or Collaborate and Promote Visibility


\bigskip

\noindent\textbf{What does ITSM focus on?}


ITSM focuses on managing IT as a set of end-to-end services that deliver value to customers and support business objectives.


\subsection{Part B – Case Analysis}

\indent 

FinServe, a mid-sized Australian financial services company, is undergoing a digital transformation. Their legacy systems are outdated, slow, and non-compliant with new financial regulations. Customers increasingly expect real-time mobile banking, faster loan approvals, and better online support.

\bigskip

To modernize operations, FinServe’s CIO has approved:

\begin{itemize}
    \item Moving 80\% of infrastructure to the cloud within 18 months.
    \item Introducing an AI-powered chatbot for customer support.
    \item Automating compliance reporting for ASIC and APRA requirements.
    \item Aligning IT services more closely with business growth strategy.
\end{itemize}

\bigskip

Challenges:

\begin{itemize}
    \item Past failures: FinServe previously tried a CRM upgrade that ran over budget and was abandoned halfway.
    \item Staff resistance: Many employees fear automation may replace jobs.
    \item Regulatory risk: Any IT failure could result in fines or reputational damage.
    \item Customer trust: Customer satisfaction has been declining; a failed rollout could lead to customer loss.
\end{itemize}

\bigskip

Frameworks to Apply:

\begin{itemize}
    \item The IT Steering Committee decides to adopt COBIT 2019 for governance and ITIL v4 for IT Service Management.
    \item They aim to avoid governance failures of the past and ensure value creation instead of just technology delivery.
\end{itemize}

\bigskip

Based on the above case study, discuss:

\noindent\textbf{How should IT governance ensure that cloud migration and AI initiatives align with business and regulatory goals?}


IT governance should establish clear strategic alignment mechanisms—linking IT initiatives to business objectives such as faster service delivery and compliance with ASIC/APRA regulations. Using COBIT’s Planning \& Organisation and Monitoring \& Evaluation domains, governance can set policies for data security, compliance checks, and performance metrics, while also involving key stakeholders (business, IT, risk, compliance teams) in decision-making.


\bigskip

\noindent\textbf{What governance failure risks might FinServe face if roles and responsibilities are not clearly defined?}


\begin{itemize}
    \item Duplication of effort or conflicting initiatives between IT and business units. 
    \item Accountability gaps where no one owns critical compliance or risk management tasks. 
    \item Decision-making delays, leading to missed deadlines or uncontrolled project scope. 
    \item Regulatory breaches, if compliance ownership is unclear, risking fines and reputational damage.
\end{itemize}


\bigskip

\noindent\textbf{Map the cloud migration project to the ITIL v4 Service Value Chain activities (Plan, Engage, Design \& Transition, Obtain/Build, Deliver \& Support, Improve).}


\begin{itemize}
    \item Plan: Define migration strategy, regulatory compliance requirements, and success criteria.
    \item Engage: Communicate with stakeholders (executives, employees, regulators, customers).
    \item Design \& Transition: Develop cloud architecture, plan data migration, and create change management plans.
    \item Obtain/Build: Acquire cloud services, configure infrastructure, and set up security controls.
    \item Deliver \& Support: Provide training, technical support, and ensure smooth operations post-migration.
    \item Improve: Gather feedback, monitor KPIs (e.g., uptime, cost savings), and refine processes continuously.
\end{itemize}


\bigskip

\noindent\textbf{What risks could emerge from introducing an AI chatbot in a financial services context? How should governance frameworks address them?}


Risks:
\begin{itemize}
    \item Incorrect or misleading financial advice to customers.
    \item Data privacy breaches if customer information is mishandled.
    \item Customer dissatisfaction due to lack of human-like empathy.
    \item Regulatory non-compliance with financial service guidelines.
\end{itemize}

Governance Responses:

\begin{itemize}
    \item Define clear boundaries of chatbot capabilities (no financial advice, only information).
    \item Implement data protection and monitoring policies.
    \item Use incident management and continual improvement practices to refine chatbot performance.
    \item Ensure accountability structures so humans oversee and validate AI decisions.
\end{itemize}


\bigskip

\noindent\textbf{If FinServe succeeds, it becomes a digital-first bank. If it fails, it risks losing customers to competitors. In your view:}


Risks:
\begin{itemize}
    \item IIs the bigger risk moving too fast (e.g., cloud migration and AI adoption without governance)
    \item or moving too slow (losing competitiveness to digital-native fintechs)?Defend your answer.
\end{itemize}
